\begin{table}[tp]\centering\scriptsize
\caption{Can a Monthly VAR Recover the Window Surprise?}\label{tab:decompdiag}
\resizebox{\ifdim\width>\textwidth\textwidth\else\width\fi}{!}{%
\begin{tabular}{lllrrrrr}\toprule Yield & FCI & Start & $R^2$, innovations (\%) & $R^2$, with 6 leads (\%) & $p$, leads & Correlation & Variance ratio\\ \midrule
Nominal & yield alone & 1973 & 2.2 & 19.0 & 0.05 & 0.11 & 99\\
Nominal & with FCI & 1973 & 0.9 & 21.2 & 0.01 & 0.02 & 75\\
Real & yield alone & 1973 & 1.4 & 17.7 & 0.05 & 0.08 & 113\\
Real & with FCI & 1973 & 0.2 & 20.0 & 0.02 & 0.01 & 77\\
Nominal & yield alone & 1996 & 3.7 & 20.4 & 0.07 & 0.19 & 54\\
Nominal & with FCI & 1996 & 1.9 & 21.6 & 0.04 & 0.11 & 35\\
Real & yield alone & 1996 & 2.6 & 17.1 & 0.21 & 0.14 & 71\\
Real & with FCI & 1996 & 1.2 & 19.5 & 0.18 & 0.07 & 44\\
\addlinespace Simulation, model as truth (median [5th, 95th]) & & & 2.5 [0.6, 5.4] & & & 0.10 [0.02, 0.21] & 65\\
\bottomrule\end{tabular}}
\begin{flushleft}\footnotesize\textit{Notes:} Monthly VAR in monthly averages (12 lags, March--December 2020 dropped). $R^2$: the month's sum of the measured surprises regressed on the month's VAR innovations (and on six months of later innovations); $p$: Wald test that the leads add nothing. Correlation: between the shock recovered from the innovations in the window direction and the measured surprise. Variance ratio: the recovered shock's variance under invertibility over the measured surprise's (1996--2026). Simulation: 200 samples generated by the real model with the FCI and the funds rate, in which the policy shock is the window surprise by construction.\end{flushleft}
\end{table}
